\chapter{Conclusiones y Líneas Futuras}

\section{Conclusiones}

\section{Líneas Futuras}

Esta línea de trabajo ha permitido obtener una visión integral del rendimiento y la eficiencia de distintos modelos de aprendizaje automático y profundo aplicados a la detección de intrusiones. Sin embargo, existen múltiples caminos de interés para continuar explorando y profundizando en esta área.

Incorporar y analizar modelos basados en secuencias temporales (como LSTM) podría ser especialmente relevante para capturar patrones de comportamiento a lo largo del tiempo, lo que es crucial en la detección de intrusiones que pueden manifestarse como anomalías temporales. 

Además, estudiar estrategias de reentrenamiento incremental permitiría evaluar cómo los modelos pueden adaptarse a nuevos datos sin necesidad de un reentrenamiento completo. Por ejemplo, se podría analizar con qué porcentaje de nuevos datos conviene reentrenar los modelos para mantener un buen rendimiento sin incurrir en un coste computacional excesivo.

Otro camino interesante sería aprovechar la similitud de características de los datasets para evaluar la transferencia entre ellos, es decir, entrenar el modelo con un dataset y probar la inferencia con otro. Esto permitiría analizar la capacidad de generalización de los modelos y su robustez.

Por último, una línea futura de gran interés sería analizar el entrenamiento para los ataques de evasión, reentrenando los modelos con muestras adversarias para mejorar su robustez frente a este tipo de amenazas, donde como se ha visto en el trabajo, los modelos pueden sufrir una degradación significativa de su eficacia.